% !TEX root = ../main.tex
\chapter{Metal and MPSGraph Integration Surface}
\label{ch:metal_mpsgraph}

The previous two chapters described the libraries that \Lucid\ uses
to talk to the GPU and the CPU: \MLX\ for GPU dispatch and
\Accelerate\ for CPU dispatch. Below both libraries is the operating
system's own graphics-and-compute substrate, \Metal, and above
\Metal\ sits another Apple-supplied tier, the \emph{Metal Performance
Shaders Graph} framework (\MPSGraph). \Lucid\ touches \Metal\
indirectly through \MLX\ and touches \MPSGraph\ directly through the
graph-capture compiler (\chref{ch:compile_design}). Both contact
points are limited in scope but consequential: the compile path is
the largest single performance feature in \Lucid\ 3.5, and the
per-op \MPSGraph\ dispatch added in 3.4
(\chref{ch:mpsgraph_dispatch}) is responsible for the most visible
single-operator speedups on small Transformer training workloads.

This chapter explains what \Metal\ and \MPSGraph\ are at a level of
detail sufficient to read the rest of the book. We do not attempt to
duplicate Apple's Metal Programming Guide; the goal is to give the
ML-framework reader a working mental model of (i)~the programming
abstractions that the operating system exposes, (ii)~the cost model
those abstractions imply, (iii)~the choice that exists between
direct \Metal\ kernels, \MPSGraph\ graphs, and the higher-level
\MLX\ runtime that sits on both, and (iv)~the specific contracts
\Lucid\ has chosen at each layer.

\section{Overview}
\label{sec:metal_mpsgraph:overview}

\Metal\ \cite{apple_metal_programming} is the operating-system-level
graphics and compute API on all of Apple's platforms (macOS, iOS,
iPadOS, tvOS, visionOS). It was introduced on iOS in 2014 and on
macOS in 2015 as a replacement for OpenGL and OpenCL. Every GPU computation on Apple silicon, no
matter what library issued it, ultimately becomes a sequence of
\Metal\ commands submitted to a \Metal\ command queue, dispatched
to the GPU through Apple's display server and the kernel-level
graphics driver.

\MPSGraph\ \cite{apple_mpsgraph} is a higher-level framework built
on top of \Metal, introduced in macOS 12 (2021). It exposes a graph IR for compute
operations and compiles graphs into multi-kernel \Metal\ executables,
with automatic kernel fusion, layout selection, and (since macOS 13)
gradient construction. \MPSGraph\ is the spiritual successor to
\code{MetalPerformanceShaders} (MPS), an earlier framework that
provided pre-canned per-op kernels without graph-level optimisation.
MPS still exists and \MPSGraph\ uses some of its primitives
internally, but for new code Apple recommends \MPSGraph.

\Lucid's relationship to these layers is layered:

\begin{itemize}
  \item \Lucid\ does not call \Metal\ directly. Every GPU command
    \Lucid\ ever issues passes through either \MLX\ (the eager
    runtime, \chref{ch:mlx_runtime}) or \MPSGraph\ (the compile
    path).
  \item \Lucid\ uses \MPSGraph\ in two places: \code{lucid.compile}
    (\chref{ch:compile_design}) emits compiled graphs onto
    \MPSGraph, and a small set of per-op fast paths in 3.4
    (\chref{ch:mpsgraph_dispatch}) dispatches selected operators
    through \MPSGraph\ even in eager mode.
  \item \Lucid\ does not implement custom Metal shaders. The few
    cases that would have benefited from one --- for example, a
    fused softmax-attention kernel --- have been deferred in favour
    of either \MLX's own optimised path or \MPSGraph\ emission.
\end{itemize}

This minimalism is deliberate. Direct Metal programming is
labour-intensive (one writes shader code in MSL, manages command
buffer lifetimes, handles synchronisation by hand), and the
ML-relevant kernels have already been written --- once in \MLX\ and
once in \MPSGraph. \Lucid's engineering budget is better spent on
the layers above.

\section{Metal Fundamentals}
\label{sec:metal_mpsgraph:metal_fundamentals}

This section describes \Metal's compute model at the level a
framework writer must understand. We omit graphics rendering (vertex
shaders, fragment shaders, render targets) entirely; \Lucid\ never
touches that path.

\subsection{The device and its command queue}
\label{ssec:metal_mpsgraph:device_queue}

The root object of every \Metal\ program is the
\code{MTLDevice}, a software handle on a single GPU. On Apple
silicon there is exactly one GPU per machine, accessible via the
function \code{MTLCreateSystemDefaultDevice()}; on Intel Macs with
discrete GPUs there could be several and the device chooser would
be relevant, but those machines are not in \Lucid's supported set.

The device is the factory for every other \Metal\ object: buffers,
textures, command queues, pipeline states. The most important child
object for compute work is the \code{MTLCommandQueue}, an ordered
FIFO into which command buffers are submitted. A single command
queue per process is the standard pattern; multiple queues are used
mainly to express concurrency between graphics and compute work.

\Lucid, via \MLX, uses exactly one command queue per process. The
queue is hidden inside \MLX's runtime and is not configurable
through \Lucid's surface.

\subsection{Command buffers and encoders}
\label{ssec:metal_mpsgraph:command_buffers}

Work is submitted to the device through \code{MTLCommandBuffer}
objects. A command buffer is built by an \emph{encoder} (one of
\code{MTLBlitCommandEncoder}, \code{MTLComputeCommandEncoder}, or
the graphics encoders) that turns API calls into device-readable
commands. For compute work the relevant encoder is
\code{MTLComputeCommandEncoder}; its primary call is
\code{dispatchThreadgroups:threadsPerThreadgroup:}, which records
a single kernel launch.

A typical compute submission is:

\begin{lstlisting}[style=lucidcpp,language=C++]
auto cmdbuf  = queue->commandBuffer();
auto encoder = cmdbuf->computeCommandEncoder();
encoder->setComputePipelineState(pso);
encoder->setBuffer(inputBuffer,  0, 0);
encoder->setBuffer(outputBuffer, 0, 1);
encoder->dispatchThreadgroups(grid, threadsPerGroup);
encoder->endEncoding();
cmdbuf->commit();
\end{lstlisting}

The command buffer is submitted (\code{commit}), and the caller may
then either continue building the next command buffer or block on
the current one with \code{waitUntilCompleted}. \MLX's lazy graph
machinery batches many primitive operations into one command buffer
to amortise the per-buffer overhead; \Lucid\ inherits this batching
without being aware of it.

\subsection{Buffers and textures}
\label{ssec:metal_mpsgraph:buffers}

Storage in \Metal\ comes in two flavours: \code{MTLBuffer} (plain
byte sequences with optional alignment requirements) and
\code{MTLTexture} (formatted, multi-dimensional images with hardware
sampling). For ML workloads only buffers matter; textures are for
graphics and image processing.

The buffer storage modes were treated in detail in
\chref{ch:unified_memory}~Section~\ref{sec:unified_memory:storage_modes}.
The summary: \MTLStorageModeShared\ for CPU/GPU shared access,
\MTLStorageModePrivate\ for GPU-only, with two other modes
(\emph{Managed}, \emph{Memoryless}) that \Lucid\ never encounters.
\MLX\ allocates all of its arrays in \MTLStorageModeShared,
which is what makes \Lucid's zero-copy bridge possible.

\subsection{Compute pipeline state}
\label{ssec:metal_mpsgraph:pso}

A \code{MTLComputePipelineState} (PSO) is a compiled, GPU-resident
representation of a single compute kernel. It is constructed from a
\code{MTLFunction} (extracted from a compiled \code{MTLLibrary}),
which is in turn compiled from MSL source code or loaded from a
\code{.metallib} binary.

The PSO is the unit of caching for compiled shader code. Constructing
a PSO from MSL source is expensive (hundreds of milliseconds for a
non-trivial kernel because it invokes the Metal shader compiler);
constructing one from a precompiled \code{.metallib} is much faster
(milliseconds, since only the device-specific final stage runs).
\MLX\ ships precompiled \code{.metallib} files inside its
distribution and constructs PSOs from them; \MPSGraph\ does the same
for the kernels its graph emitter produces.

\subsection{Synchronisation primitives}
\label{ssec:metal_mpsgraph:sync}

The relationship between two command buffers, or between a command
buffer and a CPU thread, is mediated by three \Metal\ synchronisation
primitives:

\begin{enumerate}
  \item \textbf{\code{MTLCommandBuffer::waitUntilCompleted}.}
    Blocks the calling thread until the command buffer has executed
    fully. The simplest primitive; used by \MLX's \code{eval} when
    a result is needed synchronously.
  \item \textbf{\code{MTLSharedEvent}.} A signaled event with a
    monotone 64-bit value. Both CPU threads and GPU command buffers
    can wait for the value to reach a target. This is the primitive
    that enables fine-grained CPU/GPU ordering without a full
    block; \MLX\ uses it for \code{async\_eval}.
  \item \textbf{Memory barriers within an encoder.}
    \code{MTLComputeCommandEncoder::memoryBarrier:} orders memory
    accesses within the same encoder, used to ensure that one
    kernel's writes are visible to the next kernel's reads in the
    same command buffer. Relevant only inside fused multi-kernel
    sequences; \Lucid\ inherits the correct usage from \MPSGraph.
\end{enumerate}

The choice of which primitive to use is made entirely below
\Lucid's surface. The framework's contract with the user is:
\code{Tensor} reads are synchronous; tensor operations may be
deferred but always materialise consistently by the time the result
is observed.

\section{Metal Shading Language}
\label{sec:metal_mpsgraph:msl}

The Metal Shading Language (MSL) is a C++14-derived shader language
in which Metal kernels are written. \Lucid\ does not author MSL
directly; this section is included only for the reader who wants to
understand what \MLX's kernels and \MPSGraph's emitted kernels look
like underneath.

\subsection{Kernel functions}
\label{ssec:metal_mpsgraph:kernels}

An MSL compute kernel is a function with the \code{kernel} qualifier
that takes thread-position and resource arguments:

\begin{lstlisting}[style=lucidcpp,language=C++]
#include <metal_stdlib>
using namespace metal;

kernel void add_kernel(
    device const float* a [[buffer(0)]],
    device const float* b [[buffer(1)]],
    device       float* c [[buffer(2)]],
    uint                  tid [[thread_position_in_grid]])
{
    c[tid] = a[tid] + b[tid];
}
\end{lstlisting}

The \code{[[buffer(N)]]} attribute binds the parameter to the
\code{N}-th buffer set by the encoder; \code{[[thread\_position\_in\_grid]]}
is one of several built-in thread-identity inputs. The function
runs once per thread; the encoder's
\code{dispatchThreadgroups:threadsPerThreadgroup:} determines how
many threads are launched and how they are grouped.

\subsection{Threadgroups, SIMD groups, and lanes}
\label{ssec:metal_mpsgraph:threadgroups}

\Metal's thread hierarchy has three levels:

\begin{itemize}
  \item \textbf{Grid.} The total set of threads launched by one
    dispatch. Conceptually a 1D, 2D, or 3D array; the encoder
    specifies its dimensions.
  \item \textbf{Threadgroup.} A subset of the grid that shares
    threadgroup memory and synchronises through threadgroup
    barriers. Typically 32 to 1024 threads, configured per-dispatch.
  \item \textbf{SIMD group.} A subset of a threadgroup that executes
    in lock-step on a single SIMD unit; on Apple silicon, the SIMD
    group size is 32 lanes. Threads within a SIMD group can share
    data with very low latency through SIMD-group shuffle
    intrinsics.
\end{itemize}

The mapping from these abstract levels to the silicon levels
described in \chref{ch:apple_silicon}~Section~\ref{ssec:apple_silicon:gpu_compute}
is: SIMD group $\to$ a single execution unit on a GPU core,
threadgroup $\to$ a residency on one GPU core, grid $\to$ a
dispatch across all cores. Good kernel performance requires choosing
threadgroup sizes that fit in the per-core register file
(\chref{ch:apple_silicon}~Section~\ref{ssec:apple_silicon:gpu_mem})
and using SIMD-group operations where they apply.

\subsection{Memory address spaces}
\label{ssec:metal_mpsgraph:address_spaces}

MSL distinguishes four memory address spaces:

\begin{itemize}
  \item \code{device}: pointers into \code{MTLBuffer} memory; the
    main data path.
  \item \code{threadgroup}: pointers into the threadgroup memory
    allocated for the current dispatch.
  \item \code{thread}: register-resident, per-thread local
    variables.
  \item \code{constant}: a read-only, broadcast-friendly memory
    space for small lookup tables; backed by a special hardware
    path on Apple silicon's GPU.
\end{itemize}

The address-space qualifier is mandatory on pointer types and is
checked by the MSL compiler; it gives the compiler precise
information about access patterns and enables address-specific code
generation (for instance, \code{constant} accesses are broadcast in
hardware where possible).

\subsection{Why Lucid does not write MSL}
\label{ssec:metal_mpsgraph:no_msl}

The decision to avoid hand-written MSL is the same decision, at one
level deeper, as the decision to sit on top of \MLX. The kernels
that \Lucid\ would have written (elementwise binary, reduction,
matmul, convolution, attention) have already been written by Apple
and by the \MLX\ team, with care taken about threadgroup sizing,
register pressure, SIMD utilisation, and memory access patterns.
Rewriting them in \Lucid\ would consume a kernel engineer for many
months and would not, in expectation, exceed the existing
implementations.

The escape hatch exists nonetheless. \Lucid\ exposes a
\code{lucid.metal} module
(\chref{ch:lazy_imports}) that allows the user to compile and run
an MSL kernel on a \code{Tensor}, with manual responsibility for
buffer marshalling. The interface is deliberately low-level and
unstable; it exists for research workloads that need a custom
kernel and are willing to handle Metal directly.

\section{MPSGraph}
\label{sec:metal_mpsgraph:mpsgraph}

\MPSGraph\ is the framework on which \Lucid's compile path
(\chref{ch:compile_design}) is built. Conceptually it is a graph IR
plus a compiler plus a runtime: the client constructs a directed
graph of operations on placeholder tensors, asks \MPSGraph\ to
compile the graph for a specific input-shape configuration,
receives an executable handle, and then runs the executable on
concrete input tensors.

\subsection{Graph construction}
\label{ssec:metal_mpsgraph:graph_construction}

The construction API is object-oriented Objective-C (and Swift),
with bridged calls from C++. The pattern is:

\begin{lstlisting}[style=lucidcpp,language=C++]
auto graph = MPSGraph::alloc()->init();
auto x = graph->placeholderWithShape(shape, dtype, name);
auto w = graph->placeholderWithShape(shape, dtype, name);
auto y = graph->matrixMultiplicationWithPrimaryTensor(x,
              secondaryTensor: w, name: @"matmul");
auto z = graph->additionWithPrimaryTensor(y,
              secondaryTensor: bias, name: @"bias_add");
// ... and so on
\end{lstlisting}

Every method returns an \code{MPSGraphTensor} that represents the
output of an operation; the graph is implicit in the chain of
references. There is no explicit graph-edge construction.

The operation set is large (hundreds of methods on the
\code{MPSGraph} class) and covers elementwise arithmetic,
reductions, matrix multiplication, convolution, pooling,
normalisation, activation functions, indexing, FFT, and a number of
neural-network-specific primitives (LSTM, MultiheadAttention,
SoftmaxCrossEntropy). \chref{ch:op_emitters} enumerates the subset
that \Lucid's compile path uses, which is approximately 180
primitives covering the inner loop of every supported model in the
zoo.

\subsection{Compilation and the executable cache}
\label{ssec:metal_mpsgraph:compilation}

Once constructed, a graph must be \emph{compiled} for a specific
input-shape configuration before it can run. The compilation step
is:

\begin{lstlisting}[style=lucidcpp,language=C++]
auto executable = graph->compileWithDevice(device,
    feeds: feedDict,    // placeholder -> concrete shape
    targetTensors: outputs,
    targetOperations: nil,
    compilationDescriptor: nil);
\end{lstlisting}

The returned \code{MPSGraphExecutable} is a multi-kernel
\Metal\ program ready to dispatch. Subsequent invocations on the
same shapes reuse the same executable; invocations on different
shapes require either a new compilation or, with appropriate
descriptor flags, a recompilation that the framework manages
internally.

The compilation step is expensive: hundreds of milliseconds for a
graph with a few dozen operations, into the seconds for graphs with
hundreds of operations. \Lucid's compile path therefore caches
compiled executables aggressively
(\chref{ch:mps_builder}); a training loop that compiles once and runs the same shapes for the
remainder of training amortises the compile cost to near zero.

\subsection{Backward pass}
\label{ssec:metal_mpsgraph:backward}

\MPSGraph\ in macOS 13 and later provides a
\code{gradientForPrimaryTensor:withTensors:name:} method that
constructs reverse-mode gradients directly within the graph. Given
a scalar output tensor and a list of input tensors, the method
returns a dictionary mapping each input to its gradient tensor.

\Lucid's compile path uses this method to capture forward and
backward into a single graph
(\chref{ch:fwd_bwd_compile}). The advantage is that the entire
training step --- forward, backward, gradient accumulation, parameter
update --- can be compiled as one \MPSGraph\ executable, with all
intermediate tensors fused away by the compiler. The disadvantage
is that the gradient is computed at graph-construction time, before
shape information is concrete, which constrains the operations that
admit a clean backward (and we work around the constraints with
custom \code{Function} classes for the awkward cases).

\subsection{Operation primitives}
\label{ssec:metal_mpsgraph:op_primitives}

The graph operations exposed by \MPSGraph\ fall into several
categories that map onto \Lucid's op taxonomy:

\begin{itemize}
  \item \textbf{Elementwise.} Arithmetic, comparison, logical,
    transcendental --- one \MPSGraph\ method per op.
  \item \textbf{Reduction.} Sum, mean, max, min, prod, argmax,
    argmin, with axis and keepdim parameters.
  \item \textbf{Linear algebra.} Matrix multiplication, batch
    matmul, transpose; no decompositions (Cholesky, QR, SVD are
    not in \MPSGraph\ as of macOS 26).
  \item \textbf{Convolution and pooling.} Conv1d/2d/3d (forward and
    backward), transposed convolution, max/avg pooling.
  \item \textbf{Normalisation.} Batch, layer, instance, group norm
    with learnable affine.
  \item \textbf{Indexing and reshape.} Gather, scatter, slice,
    reshape, broadcast.
  \item \textbf{Neural-network composites.} Softmax cross-entropy
    loss, LSTM, MultiheadAttention. These are higher-level than
    \MLX's primitives; \Lucid\ uses them selectively where their
    fused behaviour is preferable.
\end{itemize}

Operations that \MPSGraph\ does not provide (the most consequential
gap is the linalg family: cholesky, qr, svd, eigh) are handled in
\Lucid\ by falling back to the eager path through \MLX\ +
\Accelerate, with the result re-injected into the compiled graph as
a constant when needed. This fallback is described in
\chref{ch:compile_hardening}.

\section{MPSGraph versus MLX}
\label{sec:metal_mpsgraph:vs_mlx}

A reasonable question, given that \Lucid\ depends on both, is: when
is each the right tool? They occupy adjacent but distinct positions
in the design space; \tabref{tab:metal_mpsgraph:vs_mlx} compares
them side by side.

\begin{table}[t]
\centering
\small
\begin{tabular}{lll}
\toprule
Aspect & \MLX & \MPSGraph \\
\midrule
Abstraction level     & array library (NumPy-like) & graph IR \\
Execution model       & lazy graph, implicit eval  & explicit compile + execute \\
Per-op latency        & $\sim$10\,$\mu$s           & $\sim$5\,$\mu$s (cached) \\
First-call latency    & none                       & 0.5--3\,s (compile) \\
Op fusion             & per \texttt{eval} window   & whole-graph \\
Layout selection      & fixed                      & runtime-chosen (NCHW/NHWC) \\
Backward construction & via Lucid autograd         & via \texttt{gradientForPrimaryTensor:} \\
Dynamic shape         & natural                    & partial (3.5.1+) \\
Linalg primitives     & yes (CPU stream)           & no (Cholesky/QR/SVD absent) \\
Lucid use             & default eager path         & \texttt{lucid.compile} + 3.4 carve-out \\
\bottomrule
\end{tabular}
\caption[\MLX\ versus \MPSGraph.]{Side-by-side comparison of the
two GPU-side substrates that \Lucid\ depends on. \MLX\ is the
ergonomic eager path; \MPSGraph\ is the amortised graph-capture
path that \code{lucid.compile} targets. Both coexist within a
single \Lucid\ process because they ultimately share the same
\Metal\ command queue.}
\label{tab:metal_mpsgraph:vs_mlx}
\end{table}

\subsection{Abstraction level}
\label{ssec:metal_mpsgraph:abstraction}

\MLX\ is an array library that happens to be lazy; the graph is an
implementation detail of the runtime. The user (or, in our case,
\Lucid's dispatcher) writes operations as ordinary function calls
and reads results back when needed. The graph is constructed and
evaluated implicitly.

\MPSGraph\ is a graph library that requires explicit graph
construction, explicit compilation, and explicit execution. The user
builds the graph as a data structure, hands it to the compiler, and
runs the resulting executable. There is no eager mode.

For interactive work, model exploration, and short-running scripts,
\MLX\ is dramatically more ergonomic. For long-running training
loops in which the same operations repeat thousands of times with
the same shapes, \MPSGraph's amortised compile cost is worth the
upfront construction.

\subsection{Compilation cost}
\label{ssec:metal_mpsgraph:compile_cost}

\MLX\ has no compilation step in the same sense: the lazy graph is
constructed and evaluated per call. There is a small overhead per
operation (microseconds) for graph node construction, but no
hundreds-of-milliseconds compile latency.

\MPSGraph\ pays a compile cost the first time a graph is run with a
given input-shape configuration. For \Lucid's training-step graphs,
this is typically 0.5 to 2 seconds on first call; subsequent calls
hit the executable cache and run with the overhead of a single
command-buffer submission (microseconds). For an inference-style
workload that runs many shape variations, \MPSGraph's compile cost
dominates and \MLX\ is preferable.

\subsection{Optimisation aggressiveness}
\label{ssec:metal_mpsgraph:opt_aggression}

\MLX's lazy graph performs operator fusion and constant folding
within the lifetime of a single \code{eval} call, but does not
persist the optimised graph across calls. Each \code{eval} starts
fresh.

\MPSGraph's compiler is more aggressive: it can fuse across the
entire training step, eliminate intermediate buffers that have no
external observers, choose alternate layouts (e.g.\ NCHW vs NHWC)
based on hardware preference, and select between alternate kernel
implementations for the same operation. The training-step graphs
that \Lucid\ produces in 3.5 are typically 30 to 60\,\% smaller in
kernel count after \MPSGraph\ compilation than the equivalent
\MLX\ sequence.

\subsection{Coexistence in Lucid}
\label{ssec:metal_mpsgraph:coexistence}

\Lucid\ uses both, and the choice is explicit at the framework
boundary:

\begin{itemize}
  \item Default (eager) mode goes through \MLX. The user constructs
    a model, calls forward and backward, and observes results;
    every op is a \MLX\ primitive.
  \item \code{lucid.compile(model)} returns a \code{CompiledModule}
    that, on its first call, traces the forward pass, constructs
    an \MPSGraph, compiles, and caches. Subsequent calls run the
    cached executable directly.
  \item The 3.4 per-op \MPSGraph\ carve-out
    (\chref{ch:mpsgraph_dispatch}) is a third path: certain
    operators (GELU, Embedding backward, MaxPool backward,
    LayerNorm) dispatch through \MPSGraph\ even in eager mode,
    because their \MPSGraph\ implementations are substantially
    faster than \MLX's. This is invisible to the user but the
    dispatcher chooses it automatically.
\end{itemize}

Both paths produce the same numerical results to within a few
ULPs; the only user-visible difference is performance and the
compile-time latency of the first call.

\section{Lucid's Use of Metal and MPSGraph}
\label{sec:metal_mpsgraph:lucid_use}

Concretely, the \Lucid\ source files that touch this layer are:

\begin{itemize}
  \item \code{lucid/\_C/compile/MpsBuilder.h} and \code{.mm}: the
    Objective-C++ class that constructs \MPSGraph\ objects from
    \Lucid's trace IR. (\chref{ch:mps_builder})
  \item \code{lucid/\_C/compile/CompiledExecutable.h} and
    \code{.mm}: the wrapper around \code{MPSGraphExecutable} that
    holds compiled executables and dispatches inputs.
  \item \code{lucid/\_C/compile/ExecutableCache.h} and
    \code{.cpp}: the shape-keyed cache that maps input-shape tuples
    to compiled executables.
  \item \code{lucid/\_C/compile/OpEmitters/}: per-op emitter
    functions, one per \MPSGraph\ primitive, that turn \Lucid\ trace
    IR nodes into \code{MPSGraphTensor} calls.
  \item \code{lucid/\_C/backend/gpu/MlxBackend.cpp}: a small set of
    per-op carve-outs for the 3.4 MPSGraph dispatch.
\end{itemize}

The \code{.mm} suffix denotes Objective-C++ files, which can call
both C++ and the Objective-C \MPSGraph\ API. The bridge between
\Lucid's C++-only core and \MPSGraph's Objective-C surface is
encapsulated in these files.

Notably absent: there is no \code{MetalBackend.cpp} or
\code{MetalKernels.metal}. \Lucid\ does not own Metal kernels; the
MSL files in the build are all \MLX's. This is by design.

\section{The Compilation Cost Model}
\label{sec:metal_mpsgraph:compile_cost_model}

We now describe quantitatively the cost model that the user must
internalise for \Lucid's compile path. The numbers in this section
are taken from \chref{ch:graph_capture_wins} and reproduced here
for context.

\subsection{First-call latency}
\label{ssec:metal_mpsgraph:first_call}

The first invocation of a \code{lucid.compile(model)}-wrapped step
incurs three serial latencies:

\begin{enumerate}
  \item \textbf{Tracing.} \Lucid's tracer
    (\chref{ch:tracer}) runs the model in shape-only mode to
    record the trace IR. Typically tens to hundreds of milliseconds
    for a 100M-parameter Transformer.
  \item \textbf{Graph emission.} The MpsBuilder translates the
    trace IR into \MPSGraph\ tensor operations. Typically tens of
    milliseconds.
  \item \textbf{MPSGraph compilation.} \MPSGraph\ itself compiles
    the constructed graph into an executable. This is the dominant
    term: typically 0.5 to 3 seconds for a moderately complex
    training step.
\end{enumerate}

The total first-call latency for the GPT-2 base training step on
M3~Max is approximately 1.8 seconds. Subsequent calls hit the
executable cache and run in 110 to 115 ms.

\subsection{Shape specialisation}
\label{ssec:metal_mpsgraph:shape_spec}

The executable cache is keyed by the tuple of input shapes
(\chref{ch:compiled_module}). A training loop that varies its
batch size --- for example, last-batch-smaller-than-rest training
on a non-multiple-of-batch dataset --- will recompile per unique
shape. \Lucid\ deals with this by recommending that users drop the
last incomplete batch or pad to a fixed shape; the explicit
shape-set is part of the documented contract.

The alternative would be dynamic shape support, in which the
compiled graph accepts variable input shapes. \MPSGraph\ does
provide a partial dynamic-shape mode (with constraints on which
operations admit it), but \Lucid's 3.5 release deferred
dynamic-batch support to 3.5.1
(\chref{ch:compile_hardening}) because the parity tests on
representative models did not pass under symbolic primary shapes.

\subsection{The executable cache}
\label{ssec:metal_mpsgraph:exec_cache}

\Lucid's executable cache (\code{ExecutableCache}) has a configurable
maximum size, default 32 entries. Beyond 32, an LRU eviction policy
removes the least recently used compiled graph. For most training
workloads the working set is one or two graphs (forward+backward and
optionally an evaluation forward); the cache size is generous.

For inference workloads that compile per unique input shape, the
cache can fill quickly; users with such workloads can either raise
the cache size at module construction or pre-compile a known shape
set.

\section{Limitations and Quirks}
\label{sec:metal_mpsgraph:limits}

We close with the limitations \Lucid\ has documented from production
use of \MPSGraph.

\subsection{No backward for some forward operations}
\label{ssec:metal_mpsgraph:no_backward}

A small set of \MPSGraph\ forward operations do not have a
gradient defined. The most consequential is the linalg family
(cholesky, qr, svd, eigh), which \MPSGraph\ does not provide at all
and which therefore cannot be compiled. Some indexing operations
also have partial gradient support (scatter with non-deterministic
duplicate-index semantics is one such case).

\Lucid's compile path handles these by falling back to eager mode
for any subgraph that contains an unsupported operation
(\chref{ch:compile_hardening}). The fallback boundary is at the
op-level: one unsupported op forces eager dispatch only for that
op and its immediate consumers, not for the entire training step.

\subsection{SDK version drift}
\label{ssec:metal_mpsgraph:sdk_drift}

\MPSGraph's surface has grown each macOS release, and operations
introduced in macOS 14 are not callable on macOS 13. \Lucid's
minimum target is macOS 26.0
(\chref{ch:mlx_runtime}~Section~\ref{ssec:mlx_runtime:macos_constraint}),
which corresponds to a version of \MPSGraph\ that has the operations
\Lucid\ relies on. Backward compatibility with older macOS is not
maintained.

\subsection{Compile failure modes}
\label{ssec:metal_mpsgraph:compile_failure}

\MPSGraph\ compilation can fail for several distinct reasons:
\begin{itemize}
  \item Symbolic shapes the compiler cannot resolve.
  \item Operations not implemented for the requested dtype.
  \item Graph topology that exceeds the compiler's internal limits
    (rare but does happen for very deep, very wide graphs).
  \item Driver bugs in specific macOS point releases.
\end{itemize}

\Lucid's compile path catches the failure at compilation time and
falls back to eager mode for the entire compiled module, with a
warning logged. The user can request that the fallback be disabled
(via \code{compile(strict=True)}) to force the error to surface,
which is the right setting during model development; the default is
permissive, which is the right setting for production.

\subsection{Debugging compiled graphs}
\label{ssec:metal_mpsgraph:debug}

A compiled \MPSGraph\ executable is opaque: there is no in-process
introspection of its kernel sequence beyond the names of the
constituent operations, and no kernel-level profiler short of
Apple's \code{Instruments} application with the Metal System Trace
template. \Lucid's response to this is to dump the trace IR before
\MPSGraph\ emission (\code{lucid.compile.dump\_trace}, see
\chref{ch:tracer}), so that the user can inspect the operations
\Lucid\ requested even if the post-compilation form is invisible.

\section{Summary and Forward References}
\label{sec:metal_mpsgraph:summary}

\Metal\ is the operating-system substrate that every Apple silicon
GPU operation passes through, and \MPSGraph\ is the higher-level
graph framework that Apple supplies above it. \Lucid\ does not call
\Metal\ directly; it calls \MLX\ for eager execution and \MPSGraph\
for the compile path and a small set of per-op carve-outs. The
choice between \MLX\ and \MPSGraph\ in \Lucid\ is determined by the
user's invocation of \code{lucid.compile}: with it, the training
step is captured into an \MPSGraph\ executable for the price of a
one-time compile latency; without it, every operation flows through
\MLX's lazy graph as in any eager runtime.

This chapter closes \partref{part:hardware}. The hardware
foundation --- Apple silicon at \chref{ch:apple_silicon}, unified
memory at \chref{ch:unified_memory}, \MLX\ at \chref{ch:mlx_runtime},
\Accelerate\ at \chref{ch:accelerate}, and \Metal\ and \MPSGraph\
here --- is the complete substrate against which the rest of the
book describes \Lucid\ itself. \chref{ch:layered_dag} begins
\partref{part:system} with the layered dependency DAG that organises
the engine code, and the chapters that follow expand on each layer
of that DAG.

\begin{lucidtip}
For the user: the only \Metal-related concept that matters at the
framework boundary is \code{lucid.compile}. With it, the first
training step is several seconds longer but subsequent steps are
faster by 20--40\,\% on representative training workloads. Without
it, \MLX's eager mode runs every operation as it is called, which is
fine for development and exploration but leaves performance on the
table for production training.
\end{lucidtip}
